% T11 source: computational_method.tex lines 1--91
\label{sec:computational-method}

The computational experiments are designed to check the internal logic of M4,
identify counterexamples to overly strong interpretations, and conduct
sensitivity analysis. They do not estimate the performance of a particular team
or AI tool: the durations, topologies, and overhead functions are predominantly
synthetic. All sample sizes reported below refer either to fully enumerated
frozen matrices or to a reproducible generator, rather than to a random sample
of real-world projects.

\subsection{Experimental Questions and Reporting Rules}

Each series addresses a prespecified question and uses its own estimand.
Table~\ref{tab:experiment-design} establishes this correspondence before the
results are presented and prevents a certified optimum, the makespan of a fixed
policy, and an observation on a finite synthetic grid from being conflated.

\begin{table}[ht]
\centering
\scriptsize
\caption{Relationship between the computational series and the claims being tested}
\label{tab:experiment-design}
\begin{tabularx}{\textwidth}{@{}l>{\raggedright\arraybackslash}X>{\raggedright\arraybackslash}X>{\raggedright\arraybackslash}Xl@{}}
\toprule
Series & Question & Varied factors & Estimand and reporting rule & Method \\
\midrule
EXP-00 & Are the variants of the illustrative scenario reproduced? & Configurations 2--4 & $B_4$, $T^*$; certificate of $T^*=B_4$ & exact \\
EXP-01 & Is the bound $B_4$ always attainable? & DAG, phases, $P$ & $T^*/B_4-1$; certified positive gap & exact \\
EXP-02 & How does available capacity differ from configured parallelism? & $P$, $C(P)$, $\gamma(P)$, human-share & Profile of $T^*(P)$ & exact \\
EXP-03 & Do the aggregates identify the exact makespan? & Phase profile with equal aggregates & Paired difference in $T^*$ & exact \\
EXP-04--05 & How does granularity interact with $P$ and overhead? & Decomposition, $P$, overhead type and magnitude & Optimal granularity and interaction contrast & exact \\
EXP-06 & Are the scale and comparative conclusions robust? & Common scale and frozen error families & Invariance and preservation of ranking & exact/stress \\
EXP-07 & Do the directions of the effects carry over to synthetic DAGs? & $N$, topology, weights, placement of human-work & Paired contrasts for $T(\pi)$; separate exact audit & mixed \\
EXP-08 & What changes when a fixed $H$ is placed differently? & Concentration along a path, phases, $P$, $\gamma/C$ & Paired changes in aggregates, paths, and $T(\pi)$; exact audit & mixed \\
\bottomrule
\end{tabularx}
\end{table}

\FloatBarrier

All reports distinguish among four quantities: the structural lower bound $B_4$,
the certified optimum $T^*$, a feasible solution found without proof of
optimality, and the makespan $T(\pi)$ of a fixed policy. The notation $T^*$ is
used only for runs with status \texttt{OPTIMAL} in the computational reports.
Independently of solver status, a verified feasible schedule with
$T(\pi)=B_4$ certifies optimality through $B_4\le T^*\le T(\pi)=B_4$.
Equality must hold at the exact numerical precision of the instance;
coincidence of rounded displayed values is insufficient.

\subsection{Common Semantics and Experimental Quantities}

All series use the same M4 semantics: once a task has started, it retains one
agent slot through its final phase, including while waiting for developer
attention; all human-attention phases use a single nonpreemptive resource of
capacity 1; and a successor becomes available after all its predecessors have
completed. The multiplier $C(P)$ is applied to agent phases, whereas $\gamma(P)$
is applied to human-attention phases. The solver minimizes only the makespan;
queueing and total blocking are reported as derived characteristics but do not
serve as secondary objectives. Therefore, the reported $Q$ pertains to one
returned optimal schedule, not to the minimum $Q$ among all makespan-optimal
schedules. In the exact model, a task is assigned when its first phase begins;
the fixed simulation policy reserves a slot in advance when it selects a ready
task. Its queue and makespan are specific to that policy.

For series with an equal share of human work, define
\[
  r_h:=\frac{h_i}{k_i}.
\]
In the canonical instances of EXP-01 and EXP-03, $k_i=1$ and $r_h$ is the same
for all tasks; in the illustrative instance of EXP-02, the original $k_i$ is
retained and $h_i=r_hk_i$. The \texttt{io} profile divides $H_i$ equally around
the single agent phase: $H_i/2,A_i,H_i/2$. The \texttt{front\_loaded} profile uses
$0.9H_i,A_i,0.1H_i$. In the equal-alternation profile, the four human-attention
phases have durations $H_i/4$, and the three agent phases have durations $A_i/3$.
The historical label \texttt{alternating\_sync} denotes precisely this within-task
equality and \emph{does not} require requests from different tasks to coincide in
wall-clock time. In \texttt{alternating\_staggered}, the human-attention phases are
the same, while the agent fractions $\{1/6,2/6,3/6\}$ are cyclically permuted
across tasks. In EXP-08, the neutral name \texttt{equal\_alternating} is used for
equal alternation. Only in EXP-08 are the nominal halves, quarters, and thirds
implemented on an allocation grid of 0.001: an integer number of ticks is divided
with a remainder, which is distributed according to a seeded permutation. The
sum of the phases is preserved exactly, and parts of the same type may differ by
at most one tick.

The experimental granularity $m$ is the number of parallel subtasks that replace
the selected task. Under the canonical EXP-04--EXP-05 operator, its original
initial human-attention phase and useful agent work are divided equally among the
$m$ subtasks; all of them inherit the predecessors of the original task. They are
followed by a join that contains the original final acceptance, preserves the
original quality gate, and becomes the predecessor of all former successors. At
$m=1$, this represents semantically the same amount of useful work, although the
DAG already contains an explicit join. Therefore, the EXP-04--EXP-05 contrasts
refer to the transformed ``parts plus join'' operator; C3 in EXP-07 uses the
original task when $m=1$, and the numerical values from these two series are not
directly comparable.
